% ==================== SECTION PREAMBULE - QCM ====================

% ========== A) SYNTAXE ET PROTOTYPES ==========

\element{syntaxe}{
    \begin{question}{prototype-1}
        Quel est le prototype correct d'une fonction qui prend deux entiers en paramètres et retourne un entier ?
        \begin{multicols}{2}
            \begin{reponses}
                \bonne{\lstinline[language=C]|int calcul(int a, int b);|}
                \mauvaise{\lstinline[language=C]|calcul(int a, int b);|}
                \mauvaise{\lstinline[language=C]|int calcul(a, b);|}
                \mauvaise{\lstinline[language=C]|void calcul(int a, int b);|}
            \end{reponses}
        \end{multicols}
    \end{question}
}

% \element{syntaxe}{
%     \begin{question}{prototype-2}
%         Quelle distinction existe entre une fonction \lstinline[language=C]|void| et une fonction qui retourne un type ?
%         \begin{multicols}{2}
%             \begin{reponses}
%                 \bonne{Une fonction \lstinline[language=C]|void| ne retourne pas de valeur}
%                 \mauvaise{Une fonction \lstinline[language=C]|void| ne prend pas de parametre}
%                 \mauvaise{Une fonction \lstinline[language=C]|void| ne peut pas etre appelee}
%                 \mauvaise{Une fonction \lstinline[language=C]|void| est toujours recursive}
%             \end{reponses}
%         \end{multicols}
%     \end{question}
% }

% \element{syntaxe}{
%     \begin{question}{return-1}
%         Que se passe-t-il si on oublie le \lstinline[language=C]|return| dans une fonction qui doit retourner un entier ?
%         \begin{multicols}{2}
%             \begin{reponses}
%                 \bonne{Le comportement est indéfini, la valeur retournée est imprévisible}
%                 \mauvaise{La fonction retourne automatiquement 0}
%                 \mauvaise{Le compilateur génère une erreur qui empêche la compilation}
%                 \mauvaise{La fonction retourne NULL}
%             \end{reponses}
%         \end{multicols}
%     \end{question}
% }

% ========== B) APPELS ET PASSAGES DE PARAMETRES ==========

\element{appel}{
    \begin{question}{appel-1}
        Soit la fonction \lstinline[language=C]|float diviser(float a, float b)|. Quel appel est correct ?
        \begin{multicols}{3}
            \begin{reponses}
                \bonne{\lstinline[language=C]|float x = diviser(10.0, 2.0);|}
                \mauvaise{\lstinline[language=C]|diviser(10.0, 2.0) = x;|}
                \mauvaise{\lstinline[language=C]|float x = diviser(b:2.0, a:10.0);|}
                \mauvaise{\lstinline[language=C]|x = diviser();|}
            \end{reponses}
        \end{multicols}
    \end{question}
}

\element{appel}{
    \begin{question}{appel-2}
        Soit la fonction suivante. Qu'affichera le programme ?
        \lstinputlisting[language=C]{sources/codes/appel.c}
        \begin{multicols}{6}
            \begin{reponses}
                \bonne{12}
                \mauvaise{7}
                \mauvaise{3}
                \mauvaise{4}
                \mauvaise{0}
                \mauvaise{Erreur}
            \end{reponses}
        \end{multicols}
    \end{question}
}

\element{appel}{
    \begin{question}{ordre-param}
        Dans l'appel \lstinline[language=C]|calculer(a, b, c)|, l'ordre des parametres :
        \begin{multicols}{2}
            \begin{reponses}
                \bonne{Doit correspondre à l'ordre défini dans le prototype}
                \mauvaise{Peut être modifié librement}
                \mauvaise{N'a pas d'importance si les types sont identiques}
                \mauvaise{Est automatiquement trié par le compilateur}
            \end{reponses}
        \end{multicols}
    \end{question}
}

% ========== C) PORTEE DES VARIABLES ==========

% \element{portee}{
%     \begin{question}{portee-1}
%         Une variable déclarée à l'intérieur d'une fonction :
%         \begin{multicols}{2}
%             \begin{reponses}
%                 \bonne{Est locale et n'existe que dans cette fonction}
%                 \mauvaise{Est accessible depuis toutes les autres fonctions}
%                 \mauvaise{Est automatiquement globale}
%                 \mauvaise{Doit être obligatoirement passée en paramètre}
%             \end{reponses}
%         \end{multicols}
%     \end{question}
% }

\element{portee}{
    \begin{question}{portee-2}
        Soit le code ci-dessous. Qu'affichera le programme ?
        \lstinputlisting{sources/codes/portee-variable.c}
        \begin{multicols}{4}
            \begin{reponses}
                \bonne{Dans la fonction: 1\newline Dans le main: 10}
                \mauvaise{Dans la fonction: 11\newline Dans le main: 11}
                \mauvaise{Dans la fonction: 1\newline Dans le main: 1}
                \mauvaise{Dans la fonction: 10\newline Dans le main: 10}
            \end{reponses}
        \end{multicols}
    \end{question}
}

% \element{portee}{
%     \begin{question}{globale-vs-locale}
%         Quelle affirmation est vraie concernant les variables globales ?
%         \begin{multicols}{2}
%             \begin{reponses}
%                 \bonne{Elles sont accessibles depuis toutes les fonctions du programme}
%                 \mauvaise{Elles ne peuvent pas être modifiées}
%                 \mauvaise{Elles sont automatiquement constantes}
%                 \mauvaise{Elles sont supprimées après chaque appel de fonction}
%             \end{reponses}
%         \end{multicols}
%     \end{question}
% }

% ========== D) PREDICTIONS DE CODE ==========

\element{prediction}{
    \begin{question}{prediction-1}
        Soit le code ci-dessous. Qu'affichera le programme ?
        \lstinputlisting{sources/codes/prediction-fonction-1.c}
        \begin{multicols}{6}
            \begin{reponses}
                \bonne{13}
                \mauvaise{8}
                \mauvaise{10}
                \mauvaise{5}
                \mauvaise{16}
                \mauvaise{Erreur}
            \end{reponses}
        \end{multicols}
    \end{question}
}

\element{prediction}{
    \begin{question}{prediction-2}
        Soit le code ci-dessous. Qu'affichera le programme ?
        \lstinputlisting{sources/codes/prediction-fonction-2.c}
        \begin{multicols}{6}
            \begin{reponses}
                \bonne{10}
                \mauvaise{0}
                \mauvaise{2}
                \mauvaise{5}
                \mauvaise{8}
                \mauvaise{Erreur}
            \end{reponses}
        \end{multicols}
    \end{question}
}

\element{prediction}{
    \begin{question}{prediction-3}
        Soit le code ci-dessous. Qu'affichera le programme ?
        \lstinputlisting{sources/codes/prediction-fonction-3.c}
        \begin{multicols}{6}
            \begin{reponses}
                \bonne{7}
                \mauvaise{14}
                \mauvaise{0}
                \mauvaise{2}
                \mauvaise{21}
                \mauvaise{Erreur}
            \end{reponses}
        \end{multicols}
    \end{question}
}
